\section{Gluing at every chromatic height}\label{sec:gluing}

We now strengthen the rational estimate of \Cref{cor:rational-gluing}.
A nonzero height-\(k\) part of a cofree fixed-point algebra will produce
rational nonvanishing after adjoining \(k\) cyclic factors of arbitrarily
large order.  The rank subtraction in \Cref{lem:product-series} then
accounts for the additional \(k\) in the estimate.

\subsection{Detecting a nonzero local algebra}

\begin{lemma}[Ring detection]\label{lem:ring-detection}
Let \(R\) be a unital associative ring spectrum and \(k\ge0\).  Then
\[
 T(k)\wedge R\ne0\quad\Longrightarrow\quad K(k)\wedge R\ne0.
\]
\end{lemma}
\begin{proof}
The assertion at \(k=0\) is immediate.  For \(k>0\), choose a
positive-degree periodic self-map \(v:\Sigma^dV\to V\) of a finite
type-\(k\) spectrum whose telescope represents \(T(k)\).  Suppose that
\(K(k)\wedge R=0\).  In the associative ring spectrum
\[
 A=\End(V)\wedge R
\]
consider the element \(x=v\otimes1\in\pi_dA\).  Its image in
\(K(j)_*A\) is nilpotent at every finite height \(j\ne k\), since the
corresponding image of \(v\) is already nilpotent.  At height \(k\) it is
zero by hypothesis.  The mod-\(p\) Hurewicz image of \(v\) is nilpotent
as well: it has positive degree in the bounded graded ring
\(H_*(\End(V);\F_p)\).  This remains true after mapping to
\(H_*(A;\F_p)\).

The ring form of the nilpotence theorem therefore gives \(x^N=0\) for
some \(N\) \cite[Theorem~3(i)]{HS}.  Here the tests include rational
homology and mod-\(p\) homology, and the theorem imposes no boundedness
condition on \(A\).  Acting on \(V\wedge R\) shows that
\((v\wedge1_R)^N\) is null.  Its telescope is consequently zero, so
\(T(k)\wedge R=0\), as required.
\end{proof}

\begin{lemma}[The rational unit]\label{lem:rational-unit}
If \(k>0\) and \(B\) is a nonzero commutative algebra in
\(\Sp_{K(k)}\), then its unit has nonzero image in
\(\pi_0B\otimes\Q\).
\end{lemma}
\begin{proof}
The category \(\Sp_{K(k)}\) is stable, \(p\)-local, and
\(1\)-semiadditively symmetric monoidal
\cite[Corollary~5.3.2]{CSY}.  Apply the torsion-nilpotence theorem
\cite[Theorem~4.3.5]{CSY} with its coalgebra equal to the unit and its
commutative algebra equal to \(B\).  It says that every torsion element
of \(\pi_0B\) is nilpotent.  If the unit were torsion, it would be zero.
The unitality identity would then make \(\id_B\) zero, contradicting
\(B\ne0\).
\end{proof}

\subsection{An abelian full-mark trace}

Fix \(k>0\), and let \(E=E_k\) be Lubin--Tate theory of height \(k\)
at \(p\).  Thus \(E\) is an even periodic, \(K(k)\)-local
commutative ring with
\[
 E_0=W(\F_{p^k})[[u_1,\ldots,u_{k-1}]].
\]
We use its commutative realization from the Goerss--Hopkins--Miller
theorem \cite[Section~7]{GH}.
We write \(\widehat\otimes_E\) for tensor product in the category of
\(K(k)\)-local \(E\)-modules.  In particular, the hat denotes the
local tensor product, not an assumption of flatness on coefficients.

\begin{proposition}[Surjective characters survive base change]
\label{prop:abelian-full-mark}
Let \(B\) be a nonzero commutative \(E\)-algebra in
\(\Sp_{K(k)}\).  For every \(r\ge1\), put \(C=(C_{p^r})^k\).
Then
\[
 \bigl(U_C(B)\bigr)_{\Q}\ne0.
\]
\end{proposition}
\begin{proof}
By \Cref{lem:occurrence-fixed-points}, it suffices to show that the
full-mark projector on \(B^{BC}\) is nonzero after rationalizing
morphisms. We will detect it by computing its trace over \(E\) and
then extending scalars to \(B\).

With degree-zero coordinates for the formal
group, the complex-oriented cyclic calculation gives
\[
 E^0(BC)=E_0[[x_1,\ldots,x_k]]/
       ([p^r](x_1),\ldots,[p^r](x_k)).
\]
It is finite free of rank \(p^{rk^2}\) over \(E_0\), and
\(E^*(BC)\) is concentrated in even degrees
\cite[Corollaries~5.8 and~5.10]{HKR}.  Choosing representatives for a
degree-zero basis gives an equivalence
\[
 E^{BC}\simeq E^{\oplus p^{rk^2}}
\]
of \(E\)-modules.  Thus the categorical trace of an endomorphism of
\(E^{BC}\) is its ordinary matrix trace over \(E_0\).

Lubin--Tate theory satisfies the complete-local and finite-height
hypotheses of the generalized character theorem
\cite[Theorem~C and the examples following it]{HKR}.  After a faithfully
flat extension of \(E_0[1/p]\), the character isomorphism identifies
the base change of \(E^0(BC)\) with functions on
\(\Hom_{\cts}(\Z_p^k,C)\).  A finite \(C\)-set \(S\) acts at a
character \(\alpha\) by \(|S^{\im\alpha}|\), by the transfer formula
\cite[Theorem~D and Example~6.16]{HKR}.  Consequently the full mark
\(e_C^C\) is the projector retaining precisely the surjective
characters.  Its trace is
\begin{equation}\label{eq:full-mark-trace}
 \tr_E(e_C^C\mid E^{BC})
   =N_{k,r}\cdot1,
 \qquad
 N_{k,r}=p^{k^2(r-1)}\prod_{i=0}^{k-1}(p^k-p^i)>0
\end{equation}
in \(E_0\otimes\Q\).  Indeed every homomorphism to \(C\) factors
through \((\Z/p^r)^k\); it is onto precisely when its matrix reduces
to an invertible matrix over \(\F_p\).  Each such matrix has
\(p^{k^2(r-1)}\) lifts.  Faithful flatness descends this trace equality
from the character extension to \(E_0[1/p]\).

Now apply extension of scalars
\[
 B\widehat\otimes_E(-):
 \Mod_E(\Sp_{K(k)})\longrightarrow\Mod_B(\Sp_{K(k)}).
\]
Both module categories are \(1\)-semiadditive: their finite-groupoid
norms are computed after forgetting the module structure, where they
are equivalences.  Extension of scalars preserves colimits and hence
finite-groupoid limits \cite[Corollary~3.2.4]{CSY}.  We obtain the
canonical equivalence
\begin{equation}\label{eq:local-cochain-base-change}
 B\widehat\otimes_E E^{BC}\simeq B^{BC}.
\end{equation}
It respects the cochain algebra structures, restrictions, and
finite-cover transfers; compatibility of the canonical norms with
such functors gives the transfer assertion
\cite[Theorem~3.2.3 and Corollary~3.2.7]{CSY}.  In particular, it carries the Burnside
projector over \(E\) to that over \(B\).

This symmetric monoidal functor preserves categorical traces of
dualizable objects.  After tensoring morphism groups with \(\Q\),
\eqref{eq:full-mark-trace} therefore becomes
\[
 \tr_B(e_C^C\mid B^{BC})=N_{k,r}\cdot1_B
       \quad\text{in }\pi_0B\otimes\Q.
\]
This is nonzero by \Cref{lem:rational-unit}.  The projector is thus a
nonzero endomorphism in the rationalized module category.  It acts by
multiplication by
\(a=e_C^C\cdot1\in\pi_0(B^{BC})\otimes\Q\), so \(a\ne0\).
The rational full-mark comparison of
\Cref{lem:occurrence-fixed-points} identifies this element with the
unit of \((U_C(B))_{\Q}\), which proves the proposition.

All tensor products and groupoid limits in
\eqref{eq:local-cochain-base-change} were formed before rationalizing
morphisms.  The last step concerns the ordinary rational homotopy of
the already formed cofree spectrum.
\end{proof}

\subsection{The product comparison and the height bound}

\begin{lemma}[Coherent product comparison]\label{lem:product-comparison}
For finite groups \(P,C\) and an ordinary commutative ring spectrum
\(D\), there is a natural unit-preserving map
\[
 U_{P\times C}(D)\longrightarrow U_C(U_P(D)),
\]
where each displayed cofree construction uses the trivial action on
its coefficient algebra.
\end{lemma}
\begin{proof}
Write \(G=P\times C\). For a genuine \(G\)-spectrum \(X\), let
\(\operatorname{res}_{\mathrm{mix}}X\in\Fun(BC,\Sp^P)\) be its
restriction to a complete \(P\)-universe with trivial \(C\)-action.
Let \(\operatorname{res}_C\) take a genuine \(C\)-spectrum to its
underlying coherent \(C\)-diagram. We first justify the natural comparison
\begin{equation}\label{eq:geometric-mixed-comparison}
 \operatorname{res}_C(\Phi^P X)
   \simeq\Phi^P(\operatorname{res}_{\mathrm{mix}}X).
\end{equation}
On the left, geometric fixed points retain genuine quotient symmetry;
on the right, they are applied pointwise in \(\Fun(BC,\Sp^P)\).

Let \(\rho_P,\rho_C\) be the real regular representations, and let
\(\overline\rho_P\) be the reduced regular representation of \(P\).
Put \(V=\overline\rho_P\otimes1_C\) and
\(W=\overline\rho_P\otimes\rho_C\). For every \(H\le G\),
\[
 \dim W^H=[C:\operatorname{pr}_C H]\bigl([P:H\cap P]-1\bigr).
\]
Thus \(S^{\infty W}\) is the localization sphere for the family of
subgroups not containing \(P\). The invariant line
\(1_C\to\rho_C\) induces a unital localization comparison
\(S^{\infty V}\to S^{\infty W}\) which is a \(P\)-equivalence.
After restriction to \(\Fun(BC,\Sp^P)\), this map is an equivalence
and compares the two isotropy-separation formulas. For the right side of
\eqref{eq:geometric-mixed-comparison}, smash with \(S^{\infty V}\)
and take derived categorical \(P\)-fixed points pointwise. For the left
side, use \(S^{\infty W}\) and evaluate the normal geometric-fixed-point
formula on trivial \(C\)-representations. The comparison of spheres
therefore gives \eqref{eq:geometric-mixed-comparison}, retaining the
coherent \(C\)-action and the commutative algebra structures.

Now set \(Q=\Phi^P\Bor_G(D)\), retaining its genuine
\(C\)-structure. Model the cofree algebra as
\[
 \Bor_G(D)=F((EP\times EC)_+,\operatorname{inf}_G D).
\]
Based-space cotensors are computed levelwise, so restriction to the
chosen \(P\)-universe identifies its mixed restriction with
\(F_P((EP\times EC)_+,\operatorname{inf}_P D)\), including the
coherent \(C\)-action. Projection \(EP\times EC\to EP\) induces
a map of commutative algebras
\[
 \operatorname{const}_C\Bor_P(D)
 \longrightarrow\operatorname{res}_{\mathrm{mix}}\Bor_G(D).
\]
Forgetting \(C\), the projection is a \(P\)-equivalence: the spaces
are contractible at the trivial subgroup and have empty fixed points
at every nontrivial subgroup of \(P\). The displayed map is therefore
an equivalence in \(\Fun(BC,\Sp^P)\). Applying pointwise geometric
\(P\)-fixed points and \eqref{eq:geometric-mixed-comparison} gives
an equivalence of commutative algebras
\[
 \operatorname{res}_C Q\simeq\operatorname{const}_C U_P(D)
       \quad\text{in }\Fun(BC,\Sp).
\]
Now apply \(\Phi^C\) to the unit
\(Q\to\Bor_C(\operatorname{res}_C Q)\) of the underlying/cofree
adjunction.  Transitivity of geometric fixed points identifies its
source with \(U_{P\times C}(D)\) and its target with
\(U_C(U_P(D))\).  The construction takes place in commutative
algebras, so the map preserves units.
\end{proof}

\begin{theorem}[Height gluing]\label{thm:height-gluing}
For a finite \(p\)-group \(P\) and integers \(k,l\ge0\),
\begin{equation}\label{eq:height-gluing}
 T(k)\wedge U_P(L_{T(l)}S)\ne0
 \quad\Longrightarrow\quad k+\ell(P,1)\le l.
\end{equation}
\end{theorem}
\begin{proof}
The case \(k=0\) is \Cref{cor:rational-gluing}.  For \(l=0\),
the cofree rational calculation gives \(U_P(H\Q)=0\) for
\(P\ne1\), while \(U_1(H\Q)=H\Q\); the latter is killed by
every positive-height telescope.  Hence we may assume \(k,l>0\).

Put \(R=U_P(L_{T(l)}S)\).  This is an ordinary commutative ring
spectrum, and \Cref{lem:ring-detection} gives
\(K(k)\wedge R\ne0\).  Set
\[
 B=L_{K(k)}(E_k\wedge R).
\]
The algebra \(B\) is nonzero.  To check this without any commutativity
assumption on Morava \(K\)-theory, use the graded-field module
decomposition \cite[Proposition~1.4]{HS}
\[
 K(k)\wedge E_k\simeq\bigvee_{a\in I}\Sigma^{d_a}K(k).
\]
Here \(I\ne\varnothing\), since \(E_k\) is nonzero and
\(K(k)\)-local.  Smashing with \(R\) leaves a nonempty wedge of
nonzero spectra, so \(K(k)\wedge E_k\wedge R\ne0\).  Localization
preserves this homology.  Moreover, \(B\) is a commutative
\(E_k\)-algebra in the local symmetric monoidal category, and the
natural map \(R\to B\) is unital.

For each \(r\ge1\), take \(C=(C_{p^r})^k\).
\Cref{lem:product-comparison} and the map \(R\to B\) give unital
maps
\[
 U_{P\times C}(L_{T(l)}S)
 \longrightarrow U_C(R)\longrightarrow U_C(B).
\]
The target has a nonzero rational unit by
\Cref{prop:abelian-full-mark}.  The source unit therefore survives
rationalization as well.  Applying \Cref{cor:rational-gluing} to
\(P\times C\), we obtain
\[
 \ell(P\times(C_{p^r})^k,1)\le l\qquad(r\ge1).
\]
By \Cref{lem:product-series}, these inequalities imply
\(k+\ell(P,1)\le l\).
\end{proof}

\subsection{Finite chromatic localization}

\begin{corollary}[Finite gluing]\label{cor:finite-gluing}
For every finite \(p\)-group \(P\) and all \(m,k\ge0\),
\begin{equation}\label{eq:finite-gluing}
 m<k+\ell(P,1)\quad\Longrightarrow\quad
 T(k)\wedge U_P(L_m^fS)=0.
\end{equation}
\end{corollary}
\begin{proof}
Fix \(k\ge0\) and write \(d=\ell(P,1)\).  We induct on \(m\),
retaining all values of \(k\), including \(k>m\).  At \(m=0\),
the calculation \(U_P(H\Q)=0\) for \(P\ne1\) proves the
assertion when \(d>0\).  If \(P=1\), the required cases have
\(k>0\), and \(T(k)\wedge H\Q=0\).

Let \(m>0\) and assume \(m<k+d\).  Apply the exact functor
\(U_P\) and then smash with \(T(k)\) in the sphere fracture square
\eqref{eq:fracture}.  The corner \(U_P(L_{m-1}^fS)\) vanishes after
this test by induction.  The corner \(U_P(L_{T(m)}S)\) vanishes by
\Cref{thm:height-gluing}.

For the overlap, smashing finite localization makes
\(L_{m-1}^fL_{T(m)}S\) a module over \(L_{m-1}^fS\).
Lax monoidality therefore makes its image under \(U_P\) a unital
module over \(U_P(L_{m-1}^fS)\).  If \(T\wedge A=0\) and \(M\)
is any unital \(A\)-module, the unit and action exhibit
\[
 M\longrightarrow A\wedge M\longrightarrow M
\]
as a retraction, proving \(T\wedge M=0\).  Thus the overlap also
vanishes after smashing with \(T(k)\).  Three corners of the
resulting pullback square are zero, so its remaining corner is zero.
This completes the induction.
\end{proof}
